% ===================================================================
%  CHART No. 10 — The Sea. — The Port.
%  Traduction 2026 d'apres texto/io, controlee sur texto/fr.
%  Consigne : voir texto/en/10-table-01.tex.
%
%  Le francais decale sa numerotation d'un cran a partir de l'alinea 8
%  (il ecrit « 9. --- » sur la cle t10-08-1). C'est une coquille de
%  l'imprimeur francais, conservee dans la transcription ; l'anglais
%  suit l'ido, qui compte juste.
% ===================================================================

%%K t10-apar-1 apar
\VUsaut{4.0mm}
\VUtitre{15.0pt}{60}{CHART No. 10}
\VUsaut{3.0mm}
\VUpk{11.4pt}{40}{The Sea. --- The Port.}
\VUsaut{3.0mm}

%%K t10-01-1 p
1. --- In summer, while some look for calm and cool air in the mountains,
others prefer to go to the coast. That is what we did last year, for we
are very fond of the \VUgras{ocean} \textsuperscript{(1)}.

%%K t10-02-1 p
2. --- For my part, I take great pleasure in gazing at that immense sheet
of water stretching to the \VUgras{horizon} \textsuperscript{(2)}, and in watching
the huge \VUgras{steamships} \textsuperscript{(3)} set off on a long
\VUgras{crossing} with the travellers bound for America on board.

%%K t10-03-1 p
3. --- So my father, who knows what I like, always chooses, for the
holidays, a very quiet beach near one of our great \VUgras{naval ports}.

%%K t10-03-2 p
During our last stay, the \VUgras{squadron} came and dropped anchor behind
the huge \VUgras{breakwater} \textsuperscript{(4)} that closes and protects the
\VUgras{naval harbour} \textsuperscript{(5)}, and I was able to admire at leisure the
various \VUgras{warships}, from the small \VUgras{submarine} \textsuperscript{(10)},
which dives and disappears beneath the water, to the enormous
\VUgras{battleship} \textsuperscript{(9)}, which until now had little to fear but the
attacks of the \VUgras{torpedo boat} \textsuperscript{(8)}.

%%K t10-tit-2 sub
\VUsaut{3.0mm}
\VUpk{10.2pt}{40}{The small craft.}
\VUsaut{3.0mm}

%%K t10-04-1 p
4. --- The latter could already find, very \VUgras{cleverly}, the weak point
in the thick metal armour and plant there the dangerous
\VUgras{torpedo} whose explosion sinks the ship; but it was still less to be
feared than the submarine, for the destroyers can easily give it chase.

%%K t10-05-1 p
5. --- I was also able to see the fast armoured \VUgras{cruisers}
\textsuperscript{(6)}, almost as invulnerable as a true battleship, several
\VUgras{troopships} \textsuperscript{(12)}, three or four \VUgras{gunboats}
\textsuperscript{(7)}, and lastly a \VUgras{frigate} \textsuperscript{(11)}.
\VUgras{The sight of that fleet weighing anchor is most interesting of all.}

%%K t10-06-1 p
6. --- What a difference in build between those great masses of steel and
the elegant \VUgras{three-masters} or the graceful \VUgras{sailing ships}
\textsuperscript{(13)} still used in the merchant fleet, which I would see coming
into port under full sail as I walked along the \VUgras{pier} \textsuperscript{(14)}.
It is true that they lost much of their advantage when the \VUgras{wind} was
against them. They had to take in all sail to enter the dock, and a
\VUgras{tug} \textsuperscript{(16)} was needed to fetch them and bring them in to the
quay like plain \VUgras{barges} \textsuperscript{(17)}.

%%K t10-tit-3 sub
\VUsaut{3.0mm}
\VUpk{10.2pt}{40}{The commercial port.}
\VUsaut{3.0mm}

%%K t10-07-1 p
7. --- What makes the port I am speaking of so interesting is that it
holds not only a naval harbour, made artificially by gigantic works, but
also a marvellous \VUgras{commercial port} at the \VUgras{mouth} of a great
river.

%%K t10-08-1 p
8. --- The tongue of land separating the two docks is fortified and
defended by a line of \VUgras{batteries} and \VUgras{bastions} reaching out
into the sea. A special permit is needed to walk on the
\VUgras{ramparts} \textsuperscript{(15)}. I had obtained one easily through my
uncle, who is an engineer at the \VUgras{shipyards} \textsuperscript{(18)}, and I went
to see the ships being built under vast sheds.

%%K t10-09-1 p
9. --- I even saw a battleship launched, and I remember the thrill the
whole crowd felt when the enormous mass, left to itself, began to slide
down its cradle and came to float on the water of the river.

%%K t10-10-1 p
10. --- To reach the shipyards you cross the \VUgras{parade ground}
\textsuperscript{(19)}, where the garrison troops come to drill every day, and go
over an iron \VUgras{footbridge} \textsuperscript{(20)} linking the parade ground to
the \VUgras{arsenal} \textsuperscript{(21)}.

%%K t10-tit-4 sub
\VUsaut{3.0mm}
\VUpk{10.2pt}{40}{The Arsenal and the docks.}
\VUsaut{3.0mm}

%%K t10-11-1 p
11. --- With my uncle I visited that great establishment full of
\VUgras{guns} \textsuperscript{(22)}, \VUgras{cannonballs} \textsuperscript{(23)} and
\VUgras{shells}. I also saw the \VUgras{ironworks} \textsuperscript{(24)} where the
\VUgras{boilers} \textsuperscript{(25)} of the steamships are made.

%%K t10-12-1 p
12. --- Very near are the \VUgras{docks} \textsuperscript{(26)} --- dry docks --- and
the very remarkable \VUgras{floating docks} \textsuperscript{(27)}, in which I was
able to see close up, on a ship under repair, the \VUgras{propeller}
\textsuperscript{(28)}, the \VUgras{keel} \textsuperscript{(29)} and the \VUgras{rudder}
\textsuperscript{(30)} of a vessel.

%%K t10-13-1 p
13. --- The commercial port is quite as well worth studying as the naval
one, and I spent many hours idling on the quays where the
\VUgras{paddle steamers} \textsuperscript{(31)} that go up the river are moored, and
watching the \VUgras{sheer-legs crane} \textsuperscript{(32)} at work, lifting
enormous loads as though they were feathers.

%%K t10-tit-5 sub
\VUsaut{4.0mm}
\VUpk{10.2pt}{40}{The Pilot.}
\VUsaut{3.0mm}

%%K t10-14-1 p
14. --- I admired the \VUgras{transatlantic liners} \textsuperscript{(34)} leaving the
roadstead with their \VUgras{flags} \textsuperscript{(35)} flying, steered by a skilful
\VUgras{pilot} \textsuperscript{(36)} who knew wonderfully well how to follow the
\VUgras{channel} marked out by \VUgras{buoys} \textsuperscript{(40)} and
\VUgras{beacons}, avoiding the \VUgras{lighters} \textsuperscript{(41)} that are steered
with such difficulty by their single square sail. I could make out at the
\VUgras{bow} \textsuperscript{(37)} the second-class \VUgras{cabins} \textsuperscript{(39)},
and at the \VUgras{stern} \textsuperscript{(38)} the saloons for the first-class
passengers.

%%K t10-15-1 p
15. --- Sometimes, too, I watched the loading of a \VUgras{cargo boat}
\textsuperscript{(42)}, into which were piled the goods from the \VUgras{warehouse}
\textsuperscript{(43)} and the \VUgras{harbour station} \textsuperscript{(44)}, brought in by
the many trains of the \VUgras{quayside line} \textsuperscript{(45)}.

%%K t10-tit-6 sub
\VUsaut{3.0mm}
\VUpk{10.2pt}{40}{Rowing for pleasure.}
\VUsaut{3.0mm}

%%K t10-16-1 p
16. --- I often went boating in the \VUgras{wet dock} \textsuperscript{(53)}, where the
small \VUgras{boats} \textsuperscript{(46)} take shelter, and it was a joy to me to
handle the \VUgras{oar} \textsuperscript{(47)}. I would go up on to the \VUgras{deck}
\textsuperscript{(48)} of a \VUgras{sailing boat} \textsuperscript{(49)}, climb the
\VUgras{mast} \textsuperscript{(51)} into the \VUgras{yards} \textsuperscript{(50)} by the
\VUgras{rigging} \textsuperscript{(52)}, and then go on with my outing in a
\VUgras{skiff} \textsuperscript{(54)} among the heavy \VUgras{fishing boats}
\textsuperscript{(55)}, where the \VUgras{nets} \textsuperscript{(56)} hang drying.

%%K t10-17-1 p
17. --- On the \VUgras{quay} \textsuperscript{(58)} the most varied \VUgras{goods}
\textsuperscript{(59)} passed before me, packed in \VUgras{crates} \textsuperscript{(60)} or
in \VUgras{bales} \textsuperscript{(61)}. They made up the \VUgras{cargo} \textsuperscript{(63)}
of the barges, and powerful \VUgras{cranes} \textsuperscript{(62)} lifted them out and
set them down on \VUgras{lorries} \textsuperscript{(64)} while the \VUgras{carriers}
declared them and paid the duty at the \VUgras{customs house} \textsuperscript{(65)}.

%%K t10-18-1 p
18. --- At last, when I reached the \VUgras{swing bridge} \textsuperscript{(66)} or
the \VUgras{lock} \textsuperscript{(69)}, passing the \VUgras{dredger} \textsuperscript{(68)}
that cleans out the \VUgras{canal} \textsuperscript{(67)}, I would land on the pier
beside the \VUgras{signal mast} \textsuperscript{(70)}.

%%K t10-19-1 p
19. --- It is on the other side of that pier that the \VUgras{launches}
\textsuperscript{(71)} bringing the squadron's officers ashore come alongside. It is
a fine sight to see the sailors raise their oars in time while the boat,
carried by the \VUgras{wave} \textsuperscript{(72)}, comes easily up to the landing
steps. It is here that one can best watch the \VUgras{tide} and the movement
of the \VUgras{ebb} and \VUgras{flow} on the \VUgras{beach} \textsuperscript{(73)}.

%%K t10-tit-7 sub
\VUsaut{3.0mm}
\VUpk{10.2pt}{40}{The Officers' Club.}
\VUsaut{3.0mm}

%%K t10-20-1 p
20. --- To get home I took the \VUgras{electric tram} \textsuperscript{(74)}, whose
\VUgras{stop} \textsuperscript{(75)} is at the \VUgras{post} standing in front of the
officers' club.

%%K t10-20-2 p
From the \VUgras{balcony} \textsuperscript{(76)} of the club, decorated with
\VUgras{anchors} \textsuperscript{(77)}, guns and cannonballs, I often saw a splendid
gathering of naval officers: \VUgras{lieutenants} \textsuperscript{(78)} with their
swords, \VUgras{captains of artillery} \textsuperscript{(79)} and of
\VUgras{infantry} \textsuperscript{(80)} with their \VUgras{sabres}. Behind them stood
a \VUgras{sailor} \textsuperscript{(81)} who brought them a \VUgras{telescope} when
they wanted to make out what was happening far off.

%%K t10-21-1 p
21. --- I also saw young \VUgras{sub-lieutenants} \textsuperscript{(82)} taking their
orders from their \VUgras{commander} \textsuperscript{(83)} or from the
\VUgras{admiral} \textsuperscript{(84)}, while a \VUgras{quartermaster} \textsuperscript{(85)}
waited to carry them out to the ship.

%%K t10-22-1 p
22. --- From that balcony the view is magnificent: it commands the whole
beach with its \VUgras{tents} \textsuperscript{(86)} and \VUgras{bathing huts}
\textsuperscript{(87)}, and at bathing time one can see the \VUgras{bathers}
\textsuperscript{(88)} frolicking merrily in front of the \VUgras{casino}
\textsuperscript{(89)}.

%%K t10-23-1 p
23. --- At the end of the beach the coast forms a small rocky
\VUgras{peninsula} \textsuperscript{(91)}, where the \VUgras{breakers} \textsuperscript{(92)}
come crashing in and on which a \VUgras{lighthouse} \textsuperscript{(93)} is built,
showing sailors their way by night. This peninsula is joined to the land
only by a very narrow \VUgras{isthmus} \textsuperscript{(90)}.

%%K t10-tit-8 sub
\VUsaut{3.0mm}
\VUpk{10.2pt}{40}{A general view of the coast.}
\VUsaut{3.0mm}

%%K t10-24-1 p
24. --- The \VUgras{cliff} \textsuperscript{(94)} runs on, cut into by the sea, which
carves out of it a whimsical series of \VUgras{bays} \textsuperscript{(95)} and
\VUgras{headlands} (capes), making it most picturesque.

%%K t10-24-2 p
On the \VUgras{plateau} \textsuperscript{(96)} overlooking the \VUgras{strait}
\textsuperscript{(98)} there is a very important \VUgras{fort} \textsuperscript{(97)} and a
few villas.

%%K t10-25-1 p
25. --- That strait separates from the mainland a very dangerous
\VUgras{island} \textsuperscript{(99)}, ringed with \VUgras{rocks} \textsuperscript{(100)} and
\VUgras{reefs}, where boats and even large ships are sometimes wrecked. For
all the promptness of the rescue and the quick arrival of the
\VUgras{lifeboats}, these \VUgras{shipwrecks} are terrible, and in the
equinoctial \VUgras{gales} several ships have been lost with all hands and
cargo.

%%K t10-25-2 p
In spite of that neighbourhood, \VUgras{the port} is much used by sailors.
\VUsaut{6.0mm}
\VUfilet{22mm}
